% ---------------- 18 MODELING STRATEGY & EXPERIMENTAL ARCHITECTURE ----------------
\begin{frame}{Modeling Strategy and Hyperparameter Optimization}

\begin{block}{Optimization Framework}
Optuna's Tree-structured Parzen Estimator (TPE) was used to maximize
\textbf{Average Precision (PR-AUC)} through a Bayesian search of
\textbf{5,000 trials}.
\end{block}

\begin{itemize}
    \item End-to-end pipeline including:
    \begin{itemize}
        \item Dynamic feature engineering
        \item Scaling strategy selection (Standard, MinMax, Robust)
        \item Encoding strategy selection
        \item Model selection (Random Forest vs XGBoost)
    \end{itemize}

    \item Evaluation performed using
    \textbf{5-fold Stratified Cross Validation}.

    \item Optuna's \textbf{MedianPruner} stopped poorly performing trials early,
    reducing computational cost.

    \item Optimization target:
    \textbf{PR-AUC}, chosen because churn prediction is an imbalanced classification problem.
\end{itemize}

\end{frame}

% ---------------- 19 EXPERIMENT TRACKING BACKEND ----------------
\begin{frame}{Experiment Tracking with MLflow}

\begin{block}{Reproducible Optimization Workflow}
MLflow was integrated with Optuna to automatically track all
hyperparameter optimization experiments.
\end{block}

\begin{itemize}
    \item SQLite backend used for persistent experiment storage.
    
    \item Parent-child run structure:
    \begin{itemize}
        \item Parent run stores final model results.
        \item Child runs store individual Optuna trials.
    \end{itemize}

    \item Automatic logging of:
    \begin{itemize}
        \item Hyperparameters
        \item Cross-validation scores
        \item Test-set metrics
        \item Feature schema
        \item Serialized production pipeline
    \end{itemize}

    \item Ensures full reproducibility and auditability of model development.
\end{itemize}

\end{frame}

% ---------------- 20 EVALUATION LEADERBOARD (TABLE FORMAT) ----------------
\begin{frame}{Best Model Performance}

\centering

\small
\begin{tabular}{lc}
\hline
Metric & Value \\
\hline
Cross-Validation PR-AUC & 0.701 \\
Test PR-AUC & 0.736 \\
Accuracy & 0.865 \\
Precision & 0.781 \\
Recall & 0.598 \\
F1 Score & 0.666 \\
ROC-AUC & 0.876 \\
\hline
\end{tabular}

\vspace{0.3cm}

\begin{block}{Best Configuration}
XGBoost with Robust Scaling, 1,193 trees, learning rate 0.0267,
and class imbalance correction using
\texttt{scale\_pos\_weight = 1.67}.
\end{block}

\end{frame}

% ---------------- 21 SEARCH SPACE AUDITING DIAGNOSTICS ----------------
\begin{frame}{Hyperparameter Search Insights}

\begin{block}{Best Configuration Identified by Optuna}
\small

\begin{itemize}
    \item Model: XGBoost
    \item Scaler: RobustScaler
    \item Trees: 1,193
    \item Learning Rate: 0.0267
    \item Subsample: 0.94
    \item Column Sample Rate: 0.51
    \item Gamma: 3.33
    \item Scale Positive Weight: 1.67
\end{itemize}

\end{block}

\vspace{0.2cm}

\begin{itemize}
    \item XGBoost consistently outperformed Random Forest across the search space.
    \item Robust scaling produced the strongest results.
    \item Moderate class weighting improved minority-class detection.
    \item Bayesian optimization efficiently converged toward a stable high-performing region.
\end{itemize}

\end{frame}

% ---------------- 22 FEATURE IMPORTANCE ----------------
\begin{frame}{Key Drivers of Churn}
\begin{itemize}
    \item Low usage frequency
    \item Short customer tenure
    \item Payment delays
    \item Weak engagement score
\end{itemize}
\end{frame}

% ---------------- 23 INTERPRETABILITY ----------------
\begin{frame}{Model Explainability}
\begin{itemize}
    \item Feature importance analysis
    \item SHAP values for individual predictions
    \item Business interpretation of churn drivers
\end{itemize}
\end{frame}
